% Options for packages loaded elsewhere
% Options for packages loaded elsewhere
\PassOptionsToPackage{unicode}{hyperref}
\PassOptionsToPackage{hyphens}{url}
\PassOptionsToPackage{dvipsnames,svgnames,x11names}{xcolor}
%
\documentclass[
  letterpaper,
  DIV=11,
  numbers=noendperiod]{scrartcl}
\usepackage{xcolor}
\usepackage{amsmath,amssymb}
\setcounter{secnumdepth}{5}
\usepackage{iftex}
\ifPDFTeX
  \usepackage[T1]{fontenc}
  \usepackage[utf8]{inputenc}
  \usepackage{textcomp} % provide euro and other symbols
\else % if luatex or xetex
  \usepackage{unicode-math} % this also loads fontspec
  \defaultfontfeatures{Scale=MatchLowercase}
  \defaultfontfeatures[\rmfamily]{Ligatures=TeX,Scale=1}
\fi
\usepackage{lmodern}
\ifPDFTeX\else
  % xetex/luatex font selection
\fi
% Use upquote if available, for straight quotes in verbatim environments
\IfFileExists{upquote.sty}{\usepackage{upquote}}{}
\IfFileExists{microtype.sty}{% use microtype if available
  \usepackage[]{microtype}
  \UseMicrotypeSet[protrusion]{basicmath} % disable protrusion for tt fonts
}{}
\makeatletter
\@ifundefined{KOMAClassName}{% if non-KOMA class
  \IfFileExists{parskip.sty}{%
    \usepackage{parskip}
  }{% else
    \setlength{\parindent}{0pt}
    \setlength{\parskip}{6pt plus 2pt minus 1pt}}
}{% if KOMA class
  \KOMAoptions{parskip=half}}
\makeatother
% Make \paragraph and \subparagraph free-standing
\makeatletter
\ifx\paragraph\undefined\else
  \let\oldparagraph\paragraph
  \renewcommand{\paragraph}{
    \@ifstar
      \xxxParagraphStar
      \xxxParagraphNoStar
  }
  \newcommand{\xxxParagraphStar}[1]{\oldparagraph*{#1}\mbox{}}
  \newcommand{\xxxParagraphNoStar}[1]{\oldparagraph{#1}\mbox{}}
\fi
\ifx\subparagraph\undefined\else
  \let\oldsubparagraph\subparagraph
  \renewcommand{\subparagraph}{
    \@ifstar
      \xxxSubParagraphStar
      \xxxSubParagraphNoStar
  }
  \newcommand{\xxxSubParagraphStar}[1]{\oldsubparagraph*{#1}\mbox{}}
  \newcommand{\xxxSubParagraphNoStar}[1]{\oldsubparagraph{#1}\mbox{}}
\fi
\makeatother


\usepackage{longtable,booktabs,array}
\newcounter{none} % for unnumbered tables
\usepackage{calc} % for calculating minipage widths
% Correct order of tables after \paragraph or \subparagraph
\usepackage{etoolbox}
\makeatletter
\patchcmd\longtable{\par}{\if@noskipsec\mbox{}\fi\par}{}{}
\makeatother
% Allow footnotes in longtable head/foot
\IfFileExists{footnotehyper.sty}{\usepackage{footnotehyper}}{\usepackage{footnote}}
\makesavenoteenv{longtable}
\usepackage{graphicx}
\makeatletter
\newsavebox\pandoc@box
\newcommand*\pandocbounded[1]{% scales image to fit in text height/width
  \sbox\pandoc@box{#1}%
  \Gscale@div\@tempa{\textheight}{\dimexpr\ht\pandoc@box+\dp\pandoc@box\relax}%
  \Gscale@div\@tempb{\linewidth}{\wd\pandoc@box}%
  \ifdim\@tempb\p@<\@tempa\p@\let\@tempa\@tempb\fi% select the smaller of both
  \ifdim\@tempa\p@<\p@\scalebox{\@tempa}{\usebox\pandoc@box}%
  \else\usebox{\pandoc@box}%
  \fi%
}
% Set default figure placement to htbp
\def\fps@figure{htbp}
\makeatother


% definitions for citeproc citations
\NewDocumentCommand\citeproctext{}{}
\NewDocumentCommand\citeproc{mm}{%
  \begingroup\def\citeproctext{#2}\cite{#1}\endgroup}
\makeatletter
 % allow citations to break across lines
 \let\@cite@ofmt\@firstofone
 % avoid brackets around text for \cite:
 \def\@biblabel#1{}
 \def\@cite#1#2{{#1\if@tempswa , #2\fi}}
\makeatother
\newlength{\cslhangindent}
\setlength{\cslhangindent}{1.5em}
\newlength{\csllabelwidth}
\setlength{\csllabelwidth}{3em}
\newenvironment{CSLReferences}[2] % #1 hanging-indent, #2 entry-spacing
 {\begin{list}{}{%
  \setlength{\itemindent}{0pt}
  \setlength{\leftmargin}{0pt}
  \setlength{\parsep}{0pt}
  % turn on hanging indent if param 1 is 1
  \ifodd #1
   \setlength{\leftmargin}{\cslhangindent}
   \setlength{\itemindent}{-1\cslhangindent}
  \fi
  % set entry spacing
  \setlength{\itemsep}{#2\baselineskip}}}
 {\end{list}}
\usepackage{calc}
\newcommand{\CSLBlock}[1]{\hfill\break\parbox[t]{\linewidth}{\strut\ignorespaces#1\strut}}
\newcommand{\CSLLeftMargin}[1]{\parbox[t]{\csllabelwidth}{\strut#1\strut}}
\newcommand{\CSLRightInline}[1]{\parbox[t]{\linewidth - \csllabelwidth}{\strut#1\strut}}
\newcommand{\CSLIndent}[1]{\hspace{\cslhangindent}#1}



\setlength{\emergencystretch}{3em} % prevent overfull lines

\providecommand{\tightlist}{%
  \setlength{\itemsep}{0pt}\setlength{\parskip}{0pt}}



 


\KOMAoption{captions}{tableheading}
\makeatletter
\@ifpackageloaded{caption}{}{\usepackage{caption}}
\AtBeginDocument{%
\ifdefined\contentsname
  \renewcommand*\contentsname{Table of contents}
\else
  \newcommand\contentsname{Table of contents}
\fi
\ifdefined\listfigurename
  \renewcommand*\listfigurename{List of Figures}
\else
  \newcommand\listfigurename{List of Figures}
\fi
\ifdefined\listtablename
  \renewcommand*\listtablename{List of Tables}
\else
  \newcommand\listtablename{List of Tables}
\fi
\ifdefined\figurename
  \renewcommand*\figurename{Figure}
\else
  \newcommand\figurename{Figure}
\fi
\ifdefined\tablename
  \renewcommand*\tablename{Table}
\else
  \newcommand\tablename{Table}
\fi
}
\@ifpackageloaded{float}{}{\usepackage{float}}
\floatstyle{ruled}
\@ifundefined{c@chapter}{\newfloat{codelisting}{h}{lop}}{\newfloat{codelisting}{h}{lop}[chapter]}
\floatname{codelisting}{Listing}
\newcommand*\listoflistings{\listof{codelisting}{List of Listings}}
\makeatother
\makeatletter
\makeatother
\makeatletter
\@ifpackageloaded{caption}{}{\usepackage{caption}}
\@ifpackageloaded{subcaption}{}{\usepackage{subcaption}}
\makeatother
\usepackage{bookmark}
\IfFileExists{xurl.sty}{\usepackage{xurl}}{} % add URL line breaks if available
\urlstyle{same}
\hypersetup{
  pdftitle={A pan-Arctic database for zoobenthic biodiversity data},
  pdfauthor={Hanna Modin; Jannik Schnier; Jennifer Dannheim; Paul Renaud; Raul Primicerio; Mikael Sejr; Camille Poitrimol; Cindy Grant; Lisa Treau De Coeli; Philippe Archambault; Bodil Bluhm},
  colorlinks=true,
  linkcolor={blue},
  filecolor={Maroon},
  citecolor={Blue},
  urlcolor={Blue},
  pdfcreator={LaTeX via pandoc}}


\title{A pan-Arctic database for zoobenthic biodiversity data}
\author{Hanna Modin \and Jannik Schnier \and Jennifer Dannheim \and Paul
Renaud \and Raul Primicerio \and Mikael Sejr \and Camille
Poitrimol \and Cindy Grant \and Lisa Treau De Coeli \and Philippe
Archambault \and Bodil Bluhm}
\date{2026-03-12}
\begin{document}
\maketitle


\emph{Abstract}

As the Arctic Ocean undergoes rapid environmental change and is
subjected to increasing anthropogenic pressures, there is a growing need
for comprehensive baseline data on its biodiversity. Although sampling
efforts have been considerable over the past decades, the existing
biodiversity records are often fragmented and scattered across multiple
databases. We present BENTH-MARC, an open pan-Arctic data collection of
BENTHos in the Marine ARCtic. BENTH-MARC contains 3,614,601
georeferenced taxon records of 10,354 unique taxa from 724 datasets and
163,934 stations, spanning across the entire pan-Arctic realm and
integrating data from regional and global biodiversity databases as well
as unpublished datasets. The records in BENTH-MARC come from over two
centuries of sampling, from 1773-2025, and are based on morphological
identification. From all records, 34.4\% have associated information
about sampling gear, and 5.9\% have linked biomass data. By providing
harmonized and accessible data, BENTH-MARC can support the assessment of
biological consequences of ongoing environmental change, aid in the
identification of biodiversity hotspots and evaluation of conservation
targets. The database can also be used to guide future sampling efforts.

\section{Background \& Summary}\label{background-summary}

The Arctic Ocean is experiencing rapid environmental changes and an
increase in anthropogenic pressures\textsuperscript{1,2}, which is in
part linked to better accessibility as the sea ice
withdraws\textsuperscript{3}. The demand on Arctic resources is expected
to increase further in the future, with growing interests in among
others deep-sea mining and oil and gas extraction in the
region\textsuperscript{2}. Currently, there are knowledge gaps on how
large-scale biodiversity patterns will respond to these rapid
environmental changes and elevated anthropogenic
pressures\textsuperscript{4}. Sampling efforts in the Arctic Ocean have
been considerable over the past decades\textsuperscript{5}, and data
accessibility has increased due to improved storage options and the
growing prevalence of mandatory data publication\textsuperscript{6}.
Nevertheless, existing biodiversity records are often insufficient,
fragmented, and scattered across multiple databases\textsuperscript{7}.
This creates considerable uncertainties, posing challenges for the
sustainable management of the Arctic Ocean.

Benthic communities play a vital role in Arctic marine food webs,
nutrient recycling, and the carbon cycle\textsuperscript{8}, while also
harboring the majority of marine biodiversity\textsuperscript{9}.
Furthermore, benthic organisms are widely regarded as key indicators of
ecosystem health due to the strong coupling between environmental
drivers and benthic distribution patterns\textsuperscript{10,11}, as
well as their comparatively long lifespans. As a result, studies
investigating the impacts of environmentally hazardous activities, such
as oil and gas extraction or infrastructure development, often focus on
benthic communities. Understanding the current baseline of benthic
biodiversity is therefore critical --- not only for identifying present
biodiversity hotspots but also for predicting the future state of the
Arctic Ocean and for supporting the sustainable licensing and operation
of environmentally hazardous activities. However, sampling zoobenthos in
the Arctic Ocean presents significant challenges, including high costs,
logistical difficulties, and limited accessibility. These challenges are
particularly pronounced in deep-water areas, where research efforts
remain scarce. Most studies to date have concentrated on sea ice,
surface waters, and the water column, leaving substantial gaps in our
understanding of deep benthic environments\textsuperscript{5,12}.To
ensure that future sampling efforts target the areas of greatest need,
it is essential to create an overview of the availability, structure,
and coverage of existing data and metadata. Additionally, facilitating
reuse of already existing data helps maximize the value of each sampling
campaign.

Previous large-scale compilations of Arctic benthic biodiversity include
the works of Bluhm et al.\textsuperscript{12} and Ramirez Llodra et
al.\textsuperscript{5} on the basins, a taxon compilation by Sirenko et
al.\textsuperscript{9} on the Russian Arctic Seas, an inventory of
Arctic shelf seas by Piepenburg et al.\textsuperscript{13} and a
pan-Arctic overview of both shelves and basins by Piepenburg et
al.\textsuperscript{14}.

In a time of global change and competing land-use interests, large-scale
analyses at high taxonomic resolution are key for informing effective
science-based decision-making. Such analyses can be facilitated and
improved by access to harmonized biodiversity and biomass data. Here, we
present BENTH-MARC, an open pan-Arctic data collection of BENTHos in the
Marine ARCtic. BENTH-MARC includes 3,614,601 morphologically identified
zoobenthic records from the Global Biodiversity Information Facility
(GBIF), Ocean Biodiversity Information System (OBIS), CRITTERBASE,
Observatoire Global du Saint-Laurent (OGSL) and unpublished datasets.
The records span over the Arctic Large Marine
Ecosystems\textsuperscript{15} and date back to as early as 1773. The
records have been cleaned and harmonized and are now available for
download in different formats. By providing harmonized and accessible
data, BENTH-MARC can support the assessment of biological consequences
of ongoing environmental change and can be used to guide future sampling
efforts.

\section{Methods}\label{methods}

\subsection{Definition of the study
area}\label{definition-of-the-study-area}

The main study area consists of 18 Arctic Large Marine Ecosystems
(LMEs), which are regions of ocean space of 200,000 km2 or greater that
have unifying qualities in terms of productivity, bathymetry,
hydrography and trophically linked populations\textsuperscript{16}. The
Arctic LMEs were originally developed by the Arctic Council working
group Protection of the Arctic Marine Environment (PAME) in 2006 and
updated with revised boundaries in 2013\textsuperscript{15}. Using
ecosystem-based units supports the application of an ecosystem approach
to management, which has been promoted within the EU Marine Strategy
Framework Directive and OSPAR contexts\textsuperscript{16}.

An attempt has been made to access all zoobenthic records from within
the Arctic LMEs. Figure~\ref{fig-map_stations} shows an overview of all
sampling stations as well as the extent of the Arctic Large Marine
Ecosystems.

\begin{figure}[H]

\centering{

\pandocbounded{\includegraphics[keepaspectratio]{data_paper_manuscript_files/figure-pdf/fig-map_stations-output-1.png}}

}

\caption{\label{fig-map_stations}An overview of BENTH-MARC stations. The
Arctic Large Marine Ecosystems (purple) define the main study area. The
Large Marine Ecosystems are Labrador Sea (LS), Hudson Bay (HB), Baffin
Bay (BB), Beaufort Sea (BeS), Chukchi Sea (CS), West Bering Sea (WBS),
East Bering Sea (EBS), Aleutian Islands (AI), East Siberian Sea (ESS),
Laptev Sea (LS), Kara Sea (KS), Barents Sea (BaS), Norwegian Sea (NS),
Iceland (I), Faroe Plateau (FP), Greenland Sea (GS), North Canadian
Archipelago (NCA) and Central Arctic Ocean (CAO).}

\end{figure}%

\subsection{Data sources}\label{data-sources}

BENTH-MARC was created with the intention of including all publicly
available and morphologically identified metazoan invertebrate datasets
of meio-, macro- and megabenthos in the Arctic. All georeferenced
records belonging to the phyla Annelida, Mollusca, Arthropoda,
Echinodermata, Cnidaria, Porifera, Sipuncula, Bryozoa, Chordata,
Brachiopoda, Nematoda, Nemertea, Tardigrada, Platyhelminthes,
Hemichordata, Kinorhyncha, Priapulida, Loricifera, Phoronida and
\emph{Entoprocta} sampled from within the Arctic LMEs
(Figure~\ref{fig-map_stations}) were downloaded from GBIF and OBIS on
2026-01-31. The records were downloaded using the R packages
rGBIF\textsuperscript{17} and robis\textsuperscript{18}, which offer
programmatic access to the application programming interface (API) of
GBIF and OBIS. To access associated metadata and information,
complementary downloads targeted the related Extended Measurement or
Fact (eMoF) extensions. Data from CRITTERBASE was downloaded through
Alfred Wegener's Institute's dedicated web portal on xxxx-xx-xx. Data
from OGSL and unpublished data from Aarhus University were retrieved
through personal communication.

\subsection{Data cleaning}\label{data-cleaning}

Downloading all records in the abovementioned phyla from GBIF and OBIS
led to an inclusion of a considerable number of pelagic, terrestrial,
limnic and extinct taxa. These were at a first stage filtered out using
a comprehensive taxonomic filter at the levels of class, order, family
and species (see Supplementary material xx). As GBIF and OBIS use
different taxonomic backbones, the taxonomic filter was tailored to
correspond with the respective data source. The remaining taxa were
matched to their functional group as assigned by
WoRMS\textsuperscript{\textbf{worms?}}, which provides information on
whether the taxa can be considered benthic or pelagic. The taxa were
also matched to their respective environment through WoRMS, which
classifies all taxa as marine, brackish, fresh or terrestrial. This
enabled a further elimination of any irrelevant taxa that remained after
the first taxonomic filter. In 2022, WoRMS had 81\% coverage in terms of
functional group and 100\% coverage in terms of environment for
accepted, extant, marine species\textsuperscript{\textbf{worms2?}}. Only
records belonging to extant marine benthic taxa were kept for further
processing. All information about functional groups, environment and
status as extant or extinct was accessed through the R package
worrms\textsuperscript{\textbf{worrms?}}.

All records belonging to datasets with molecular-, telemetry- or
acoustic-based identification or with a main focus on zooplankton, fish
or sympagic taxa were removed. Citizen-science datasets without
expert-level quality control have been deleted. Records sampled with
typically pelagic gear such as plankton net and continuous plankton
recorder were excluded and fossil specimens were filtered out. Our
filters also targeted the georeferencing quality, removing all records
with coordinate uncertainty of \textgreater0.01 degrees or
\textgreater10,000 meters. Records with an unspecified geographical
uncertainty were kept. Records with the known default uncertainty values
301, 3036, 999, and 9999 were removed, as they typically have a higher
uncertainty than stated by the given value. This was done in accordance
with GBIF recommendations\textsuperscript{19}. All records denoting
absence were excluded.

CRITTERBASE, GBIF and OBIS have considerable data overlap. CRITTERBASE
was considered the most complete source in terms of metadata and
therefore preferred over the other two sources when possible. GBIF has
more easily accessible dataset information than OBIS and was therefore
given secondary priority. An effort was made to add metadata from OBIS
records to GBIF records, whenever there were duplicates between these
databases and more metadata had been stored in OBIS.

\subsection{Data compilation}\label{data-compilation}

AphiaIDs were retrieved for all taxa in the GBIF subset using the r
package worrms\textsuperscript{\textbf{worrms?}}, using the first match
and manually processing the taxa where no match could be found. The
other sources already had AphiaIDs assigned to each taxon. To ensure
that all records were updated against the latest taxonomic information
available on WoRMS, all records from all sources were stripped of all
taxonomic information apart from AphiaID, including scientific name,
rank, classification and taxon status. The taxonomic information was
then harvested anew from WoRMS using
worrms\textsuperscript{\textbf{worrms?}}, to ensure that all taxonomic
information is up-to-date and has been retrieved on the same date for
all taxa.

\section{Data records}\label{data-records}

The data is primarily available as two separate files; one Parquet file
with all occurrence records in a long format, and one CSV file with
dataset-level metadata. The column datasetKey links these two files.
Parquet is a column-oriented data storage format that facilitates
efficient data processing and retrieval of big data. Unlike row-based
formats like CSV, Parquet organizes data by columns, which allows for
better compression and faster analytical queries\textsuperscript{20}. To
accommodate users who cannot access the data in Parquet format, the
occurrence records are also available as one CSV file per LME.

All files are available on Zenodo/other website.

\section{Data overview}\label{data-overview}

BENTH-MARC contains 3,614,601 georeferenced taxon records of 10,354
unique taxa from 724 datasets and 163,808 stations. All unique AphiaIDs
in the database are here seen as unique taxa, even if the taxon has not
been identified to species level. This means that the same taxon may be
identified at both species level and higher taxonomic levels, and the
number of unique taxa is therefore an overestimation of the number of
unique species. The records in BENTH-MARC represent two and a half
centuries of sampling, spanning from 1773 to 2025.
Figure~\ref{fig-tax_resolution} shows the taxonomic resolution of the
records. The majority of the records (61 \%) has been identified to
species level, followed by genus (11 \%) and family level (4 \%).

\begin{figure}[H]

\centering{

\pandocbounded{\includegraphics[keepaspectratio]{data_paper_manuscript_files/figure-pdf/fig-tax_resolution-output-1.png}}

}

\caption{\label{fig-tax_resolution}Taxonomic resolution of all records
of benthic invertebrates in the BENTH-MARC data collection. The bars
show records per rank from phylum to subspecies.}

\end{figure}%

Figure~\ref{fig-top10phyla} shows an overview of the top ten phyla in
BENTH-MARC in terms of number of records. The number of records can be
compared to the number of unique taxa per phylum. Annelida is by far the
phylum with the highest number of records, with a share of 46 \% of the
total number of records. However, the highest number of unique taxa
(2,536) can be found within the phylum Arthropoda, even though it only
has 14 \% of the total number of records. The phylum Bryozoa stands out
in terms of having the highest taxa to records ratio.

\begin{figure}[H]

\centering{

\pandocbounded{\includegraphics[keepaspectratio]{data_paper_manuscript_files/figure-pdf/fig-top10phyla-output-1.png}}

}

\caption{\label{fig-top10phyla}The number of records and taxa of the ten
most records-rich phyla of benthic invertebrates in the BENTH-MARC data
collection, split into species-level (taxonRank = SPECIES) and other
taxonomic levels.}

\end{figure}%

The records in BENTH-MARC have varying degrees of associated information
about biomass, occurrence metrics and sampling gear
(Figure~\ref{fig-panelmap} a-c). We have considered only records with
information about both number of individuals and sampling area to be
true abundance records, as no standardization methods can be applied on
records where only one of these is known. Stations with full biomass
coverage and with abundance data are mainly located in the Barents Sea,
the Norwegian Sea, Chukchi Sea and Laptev Sea. In terms of the number of
records, only 5.9\% have biomass data and 6.2\% have abundance data.
Since stations can have a combination of records with and without
biomass coverage as well as a combination of records with different
types of occurrence metrics, the categories partial (biomass coverage)
and mixed (occurrence metrics) are displayed in the maps. In addition,
34.4\% of the records have associated information about sampling gear.

As seen in Figure~\ref{fig-panelmap} d, most stations have low numbers
of records. 73 \% have less than ten records per station, and 90\% have
less than 64 records per station. The lowest number of records per
station is one and the highest number is 295.

\pandocbounded{\includegraphics[keepaspectratio]{data_paper_manuscript_files/figure-pdf/cell-25-output-1.png}}

\begin{figure}[H]

\centering{

\pandocbounded{\includegraphics[keepaspectratio]{data_paper_manuscript_files/figure-pdf/fig-panelmap-output-1.png}}

}

\caption{\label{fig-panelmap}a) Biomass coverage, b) occurrence metrics
and c) gear information across stations and records, as well as d)
number of unique taxa per station of benthic invertebrates in the
BENTH-MARC data collection. In 4 d), stations with higher numbers of
taxa are drawn on top of stations with lower numbers of taxa.}

\end{figure}%

\section{Technical validation}\label{technical-validation}

\section{Usage notes}\label{usage-notes}

The data in BENTH-MARC come with certain qualities that need to be
accounted for when applying the data on larger-scale analyses. While
Figure~\ref{fig-panelmap} can help inform the user on the suitable
selection of data for a certain scope, there is also an inherent
patchiness in the data that the user needs to be aware of. The data are
unevenly distributed over regions, habitat types and depths and the
sampling focus has shifted over time. These characteristics are
described below. Furthermore, varying levels of taxonomic resolution in
the data can lead to a need for normalizing the records prior to making
comparisons (Figure 2).

While the spatial distribution of stations across Arctic LMEs can be
seen in Figure~\ref{fig-map_stations}, Figure~\ref{fig-lme} shows an
overview of the number of records and unique taxa per LME. Sampling
effort in terms of number of records varies considerably between LMEs,
with 84.6 \% of the total number of records having been collected in the
two most well-sampled LMEs; the Norwegian Sea and Barents Sea. The
availability of public records is significantly lower in the remaining
LMEs, and the five most records-poor LMEs East Siberian Sea, Northern
Canadian Archipelago, Fareo Plateau, West Bering Sea, and Iceland
contribute with merely 0.6 \% of the total number of records. The
Norwegian Sea is the most well-sampled LME, while Barents Sea has higher
numbers of unique taxa. Iceland stands out in terms of having a low
number of records but a disproportionally high number of unique taxa.

\begin{figure}[H]

\centering{

\pandocbounded{\includegraphics[keepaspectratio]{data_paper_manuscript_files/figure-pdf/fig-lme-output-1.png}}

}

\caption{\label{fig-lme}Number of records and unique taxa of benthic
invertebrates in the BENTH-MARC data collection per Arctic Large Marine
Ecosystem. The colors illustrate the proportion of species-level data
compared with data on other taxonomic levels.}

\end{figure}%

Similarly, sampling effort varies considerably over geomorphological
features on a pan-Arctic scale. In Figure~\ref{fig-geo}, the global
geomorphological features mapped by Harris et al.~(2014) have been
merged into a coarser set of categories to visualize the distribution of
sampling effort in terms of number of records as well as the number of
unique taxa per category. As seen in Figure~\ref{fig-geo}, the deep and
medium shelves are overrepresented in terms of both number of records.
This reflects the high sampling pressure in Barents Sea and the
Norwegian Sea, which are both rich in fjords that are deep enough to be
categorized as medium and deep shelf.

The shelves are most easily accessible and the deep, medium and shallow
shelf have together received a staggering 88.7 \% of the overall
sampling effort in terms of number of records sampled. At the other end
of the spectrum, the hadal zones together with abyssal plains, mountains
and hills have collectively only received 1.2 \% of the total sampling
effort. Due to near-shore locations that fall outside the coverage of
the shapefiles provided by Harris et al.~(2014), 3.9 \% of the records
are categorized as Outside geomorphological categories.

The proportional relationship between number of records and number of
unique taxa changes with sampling pressure and depth. The
taxa-to-records ratio is lowest at the deep shelf at 0.47 and generally
gradually increases as the number of records decreases and the depth
increases, with abyssal plains and hadal zones having the highest ratio
at 16.24. There are a few exceptions to this general pattern, such as
the shallow and medium shelf having fewer records and a higher ratio
than the deep shelf, and ridges having fewer records and a lower ratio
than the abyssal mountains and hills. The deeper areas are generally
less well-sampled and each sample from these areas therefore has a
higher probability of encountering taxa that had not been previously
recorded in the area, thereby explaining the higher taxa-to-records
ratios. The ratios indicate that the low numbers of unique taxa in the
deep areas are a consequence of low sampling effort rather than a true
representation of the actual biodiversity in these areas, and that the
deep areas are likely to be considerably more diverse than current data
suggests.

\pandocbounded{\includegraphics[keepaspectratio]{data_paper_manuscript_files/figure-pdf/cell-36-output-1.png}}

\begin{figure}[H]

\centering{

\pandocbounded{\includegraphics[keepaspectratio]{data_paper_manuscript_files/figure-pdf/fig-geo-output-1.png}}

}

\caption{\label{fig-geo}An overview of the number of records and unique
taxa of benthic invertebrates in the BENTH-MARC database per
geomorphological category. The geomorphological categories have been
simplified from Harris et al.~(2014).}

\end{figure}%

As seen in Figure~\ref{fig-depth}, the vast majority of sampling has
occurred at shallow depths. 50 \% of the records have been sampled at a
depth of less than 180 meters, and only 29 \% have been sampled at a
depth greater than 300 meters. 98 \% of the records have been sampled at
the top 1000 meters.

\begin{figure}[H]

\centering{

\pandocbounded{\includegraphics[keepaspectratio]{data_paper_manuscript_files/figure-pdf/fig-depth-output-1.png}}

}

\caption{\label{fig-depth}Number of records of benthic invertebrates in
the BENTH-MARC database over depth. Note the different scales in the
split x axis.}

\end{figure}%

The data in BENTH-MARC span back all the way to the 1700s, but only 0.7
\% of the records originate from before the 20th century. The vast
majority of the data are recent, with 88.3 \% of the records being
sampled after year 2000 (Figure~\ref{fig-time}). Apart from the
increasing number of records over time, there also appears to have been
a shift in focus from sampling only a few taxonomic groups at a time to
more community-oriented sampling. Figure~\ref{fig-time} a) shows the
development over time from almost exclusively single-phylum stations to
the currently lower proportion of stations that only focus on one
phylum. Apart for over the past three years, there has been a clear
downward trend since around year 2000. The relative increase over the
past three years might be related to a time lag in the identification of
taxa in more comprehensive, community-based studies. The overall
reduction in single-phylum stations can be interpreted as a shift in
research focus from taxonomy to ecology. Together, the limited
availability of historic data and the shift in study focus limit the
applicability of the data in long-term comparisons.

Figure~\ref{fig-time} b) shows the six most common phyla among
single-phylum stations. Arthropoda, Molllusca, Cnidaria, Porifera and
Echinodermata have been the focus of more than 10,000 stations, whereas
Annelida has been the focus of fewer targeted studies.

\pandocbounded{\includegraphics[keepaspectratio]{data_paper_manuscript_files/figure-pdf/cell-43-output-1.png}}

\begin{figure}[H]

\centering{

\pandocbounded{\includegraphics[keepaspectratio]{data_paper_manuscript_files/figure-pdf/fig-time-output-1.png}}

}

\caption{\label{fig-time}a) shows the number of records over time and
the proportion of single-phylum stations per year, indicating a shift in
sampling focus from single-phylum stations to communities. Figure 8: b)
shows the phyla that have most commonly been the focus of single-phylum
stations.}

\end{figure}%

\section{Data availability}\label{data-availability}

All data described in this paper can be accessed as Parquet or CSV files
on Zenodo/other website. Usage licence is under Creative Commons
Attribution 4.0 International (CC-BY-4.0).

\section{Code availability}\label{code-availability}

The R and Python code that was used to create the database BENTH-MARC
and to generate figures 1-8 can be found in the GitHub repository
arctic-benthic-biodiverity\textsuperscript{21}.

\section{Author contributions}\label{author-contributions}

The first authorship is shared between H.M. and J.S., who performed the
majority of the data curation. In addition, J.D. also contributed with
data curation. M.S., C.P., C.G., L.T.D.C., and P.A. contributed with
unpublished datasets. H.M., B.B., J.D., P.R., R.P., and M.S. developed
the conceptual design of the paper. H.M. wrote the manuscript draft and
all other authors contributed with manuscript revisions.

\section{Competing interests}\label{competing-interests}

No competing interests are declared for this work.

\section{Funding}\label{funding}

This work was funded by the Horizon Europe project POMP - Polar Ocean
Mitigation Potential (grant agreement 101136875). The POMP project has
been approved under HORIZON-CL6-2023-CLIMATE-01-3: Ocean and coastal
waters carbon- and biodiversity-rich ecosystems and habitats in Europe
and the Polar Regions.

\section*{References}\label{references}
\addcontentsline{toc}{section}{References}

\protect\phantomsection\label{refs}
\begin{CSLReferences}{0}{0}
\bibitem[\citeproctext]{ref-overland2019}
\CSLLeftMargin{1. }%
\CSLRightInline{Overland, J. \emph{et al.}
\href{https://doi.org/10.1016/j.polar.2018.11.008}{The urgency of arctic
change}. \emph{Polar Science} \textbf{21}, 6--13 (2019).}

\bibitem[\citeproctext]{ref-bidgood2024}
\CSLLeftMargin{2. }%
\CSLRightInline{Bidgood, A. K. \& Hall, J.
\href{https://doi.org/doi:10.3389/esss.2024.10117}{We need to talk about
mining in the arctic}. \emph{Earth Science, Systems and Society}
\textbf{4}, 10117 (2024).}

\bibitem[\citeproctext]{ref-pame2025}
\CSLLeftMargin{3. }%
\CSLRightInline{PAME. \emph{The Increase in Arctic Shipping: 2013-2024}.
\url{https://hdl.handle.net/11374/2733.3} (2025).}

\bibitem[\citeproctext]{ref-brandt2023}
\CSLLeftMargin{4. }%
\CSLRightInline{Brandt, S., Wassmann, P. \& Piepenburg, D.
\href{https://doi.org/10.3389/fmars.2023.1096222}{Revisiting the
footprints of climate change in arctic marine food webs: An assessment
of knowledge gained since 2010}. \emph{Frontiers in Marine Science}
\textbf{10}, (2023).}

\bibitem[\citeproctext]{ref-ramirezllodra2024}
\CSLLeftMargin{5. }%
\CSLRightInline{Ramirez-Llodra, E. \emph{et al.}
\href{https://doi.org/10.1525/elementa.2023.00140}{The emerging picture
of a diverse deep arctic ocean seafloor: From habitats to ecosystems}.
\emph{Elementa-Science of the Anthropocene} \textbf{12}, (2024).}

\bibitem[\citeproctext]{ref-michner2015}
\CSLLeftMargin{6. }%
\CSLRightInline{Michener, W. K.
\href{https://doi.org/10.1016/j.ecoinf.2015.06.010}{Ecological data
sharing}. \emph{Ecological Informatics} \textbf{29}, 33--44 (2015).}

\bibitem[\citeproctext]{ref-caff2017}
\CSLLeftMargin{7. }%
\CSLRightInline{CAFF. \emph{State of the Arctic Marine Biodiversity
Report}. \url{http://hdl.handle.net/11374/1945} (2017).}

\bibitem[\citeproctext]{ref-kedra2015}
\CSLLeftMargin{8. }%
\CSLRightInline{Kędra, M. \emph{et al.}
\href{https://doi.org/10.3402/polar.v34.23775}{Status and trends in the
structure of arctic benthic food webs}. \emph{Polar Research}
\textbf{34}, 23775 (2015).}

\bibitem[\citeproctext]{ref-sirenko2001}
\CSLLeftMargin{9. }%
\CSLRightInline{Sirenko, B. \emph{et al.} \emph{List of Species of
Free-Living Invertebrates of Eurasian Arctic Seas and Adjacent Deep
Waters}. (2021).}

\bibitem[\citeproctext]{ref-kraan2010}
\CSLLeftMargin{10. }%
\CSLRightInline{Kraan, C., Aarts, G., Meer, J. van der \& Piersma, T.
\href{https://doi.org/Doi\%2010.1890/09-2040.1}{The role of
environmental variables in structuring landscape-scale species
distributions in seafloor habitats}. \emph{Ecology} \textbf{91},
1583--1590 (2010).}

\bibitem[\citeproctext]{ref-kraan2015}
\CSLLeftMargin{11. }%
\CSLRightInline{Kraan, C., Dormann, C. F., Greenfield, B. L. \& Thrush,
S. F.
\href{https://doi.org/ARTN\%20e0142411\%0A10.1371/journal.pone.0142411}{Cross-scale
variation in biodiversity-environment links illustrated by coastal
sandflat communities}. \emph{Plos One} \textbf{10}, (2015).}

\bibitem[\citeproctext]{ref-bluhm2011}
\CSLLeftMargin{12. }%
\CSLRightInline{Bluhm, B. A. \emph{et al.}
\href{https://doi.org/10.1007/s12526-010-0078-4}{Diversity of the arctic
deep-sea benthos}. \emph{Marine Biodiversity} \textbf{41}, 87--107
(2011).}

\bibitem[\citeproctext]{ref-piepenburg2011}
\CSLLeftMargin{13. }%
\CSLRightInline{Piepenburg, D. \emph{et al.}
\href{https://doi.org/10.1007/s12526-010-0059-7}{Towards a pan-arctic
inventory of the species diversity of the macro- and megabenthic fauna
of the arctic shelf seas}. \emph{Marine Biodiversity} \textbf{41},
51--70 (2011).}

\bibitem[\citeproctext]{ref-piepenburg2024}
\CSLLeftMargin{14. }%
\CSLRightInline{Piepenburg, D. \emph{et al.}
\href{https://doi.org/10.5194/essd-16-1177-2024}{PANABIO: A
point-referenced PAN-arctic data collection of benthic BIOtas}.
\emph{Earth System Science Data} \textbf{16}, 1177--1184 (2024).}

\bibitem[\citeproctext]{ref-pame2013}
\CSLLeftMargin{15. }%
\CSLRightInline{PAME. \emph{Large Marine Ecosystems (LMEs) of the Arctic
Area - Revision of the Arctic LME Map}.
\url{http://hdl.handle.net/11374/61} (2013).}

\bibitem[\citeproctext]{ref-pamend}
\CSLLeftMargin{16. }%
\CSLRightInline{PAME.
\href{https://pame.is/ourwork/ecosystem-approach-to-management-ea/large-marine-ecosystems/}{Large
marine ecosystems}. vol. 2025.}

\bibitem[\citeproctext]{ref-rgbif2026}
\CSLLeftMargin{17. }%
\CSLRightInline{S., C. \emph{et al.}
\href{https://CRAN.R-project.org/package=rgbif}{Rgbif: Interface to the
global biodiversity information facility API}. (2026).}

\bibitem[\citeproctext]{ref-robis2022}
\CSLLeftMargin{18. }%
\CSLRightInline{Provoost, P. \& Bosch, S.
\href{https://CRAN.R-project.org/package=robis}{Robis: Ocean
biodiversity information system (OBIS) client}. (2022).}

\bibitem[\citeproctext]{ref-gbif2021}
\CSLLeftMargin{19. }%
\CSLRightInline{Waller, J.
\href{https://data-blog.gbif.org/post/gbif-filtering-guide/}{Common
things to look out for when post-processing GBIF downloads}. (2021).}

\bibitem[\citeproctext]{ref-parquet2022}
\CSLLeftMargin{20. }%
\CSLRightInline{Parquet, A.
\href{https://parquet.apache.org/docs/overview/motivation/}{Motivation}.
(2022).}

\bibitem[\citeproctext]{ref-githubrep}
\CSLLeftMargin{21. }%
\CSLRightInline{Modin, H.
\href{https://github.com/hanna-modin/arctic-benthic-biodiversity}{Arctic-benthic-biodiversity}.
vol. 2026 (2026).}

\end{CSLReferences}




\end{document}
